% Generated draft: requires mathematical and source review.
\documentclass[11pt]{article}
\usepackage[margin=25mm]{geometry}
\usepackage{fontspec}
\setmainfont{Latin Modern Roman}
\usepackage{amsmath,amssymb,mathtools,unicode-math}
\setmathfont{Latin Modern Math}
\usepackage{enumitem,needspace,xurl,hyperref}
\setlength{\parindent}{0pt}
\setlength{\parskip}{5pt}
\setlength{\emergencystretch}{3em}
\newcommand{\N}{\mathbb{N}}
\newcommand{\Z}{\mathbb{Z}}
\newcommand{\Q}{\mathbb{Q}}
\newcommand{\R}{\mathbb{R}}
\newcommand{\C}{\mathbb{C}}
\newcommand{\F}{\mathbb{F}}
\newcommand{\A}{\mathbb{A}}
\newcommand{\PP}{\mathbb{P}}
\newcommand{\E}{\mathbb{E}}
\newcommand{\eps}{\varepsilon}
\newcommand{\dd}{\,\mathrm{d}}
\begin{document}
\section*{UnsolvedMath ID: 800339}
\textbf{Source problem code:} catalogue entry
\subsection{Definitions and mathematical statement}
A domino is a pair of unit lattice squares sharing a full edge. For a finite region \(R\) consisting of unit lattice squares, let \(\mathcal{T}(R)\) be the set of partitions of its squares into dominoes, and define
\[
T(R)=\#\mathcal{T}(R).
\]
Let \(I\) denote the intended range of the parameter \(n\). For \(n\in I\), let \(R_n\) denote the region called a \((2n-1)\)-by-\(2n\) Aztec rectangle in the supplied source, and let \(c_n\) denote its central square. Write
\[
A_n=\{s\in R_n:s\text{ shares a full edge with }c_n\}.
\]
The source describes \(A_n\) as consisting of four squares. The supplied excerpt does not specify the geometric convention defining this Aztec rectangle, identify its central square by coordinates, or give the intended parameter range \(I\). Thus the following formulation is conditional on recovering those source conventions, including any restrictions needed for all four adjacent squares to exist.

\textbf{Enumeration problem (Propp, 1998, Problem 10).} Determine an explicit formula for
\[
N(n,s)=T\bigl(R_n\setminus\{s\}\bigr)
\qquad\text{for every }n\in I\text{ and every }s\in A_n.
\]
Here deletion means removal of the indicated unit square. The supplied source reports that
\[
T\bigl(R_n\setminus\{c_n\}\bigr)=0,
\qquad
T\bigl(R_n\setminus\{s\}\bigr)>0\quad(s\in A_n),
\]
and describes the latter counts as appearing to be a \emph{nice round number}. These are supplied source assertions and a qualitative observation, respectively; they are not independently verified here. The excerpt does not explicitly assert that the four choices of \(s\) have equal counts.
\subsection{Short English statement}
Propp asks for a formula counting domino tilings after deleting any square sharing an edge with the central square of a \((2n-1)\)-by-\(2n\) Aztec rectangle. The exact Aztec rectangle convention and admissible values of \(n\) must be recovered from the source. The registry labels the problem open; this is a supplied status, not an independently verified assessment.
\subsection{Sources}
\begin{itemize}
\item Local UnsolvedMath registry, permanent ID \(800339\), exact source code \texttt{catalogue entry}, title \emph{Enumeration of Matchings --- Problem 10}. The supplied record attributes the problem to Propp, with proposed year 1998, and labels its status \texttt{open}.
\item Propp, \emph{Twenty Open Problems in Enumeration of Matchings} (1998), Problem 10 (\#10), \url{https://arxiv.org/abs/math/9801060}. The supplied provenance records access on 2026-10-16 and extraction by \texttt{arxiv\_latex}; these provenance claims are supplied and were not independently verified here.
\end{itemize}
\end{document}
